\section{Descent through a locally constant MRC fibration}
\label{V:sec-mrc-descent}

We apply Theorem~\ref{V:controlled-invariant-nonvanishing} to the
rationally connected fibers of a locally constant fibration.  The constant extension of a section
on one fiber gives a relative anticanonical section precisely when
the section is fixed by monodromy.  The argument below first reduces monodromy to
a connected commutative algebraic group, then uses its unipotent part
to pass from torus invariance to monodromy invariance.  The final two
steps remove the base canonical bundle and descend through a finite
cover.  This is the invariant-section reduction of
M\"uller~\cite[Theorem~2.2 and its proof]{Muller2025}, with the
fiberwise section supplied here by the controlled-cone theorem.

Let $X$ be a smooth connected projective complex variety, and assume
that $-K_X$ has a smooth Hermitian metric with semipositive curvature.
The structure input is Matsumura--Wang's theorem for projective klt
pairs with nef anti-log-canonical divisor
\cite[Theorem~1.1, Definition~2.3, Theorem~4.12, and the conclusion
of~\S4.4]{MatsumuraWangV3}.  We use its zero-boundary case.  The
hypotheses hold because smoothness implies the klt condition and
smooth semipositivity implies nefness.  The theorem gives a finite
quasi-\'etale cover $p_0:X_0\to X$ and a holomorphic maximal
rationally connected fibration
\[
 f_0:X_0\longrightarrow B_0,
 \qquad K_{B_0}\sim_{\mathbb Q}0,
\]
where $B_0$ is projective and the fibration is locally constant.
Concretely, for a fixed connected rationally connected fiber $F$, there
are a representation
\[
 \rho_0:\pi_1(B_0)\longrightarrow\operatorname{Aut}(F)
 \quad\text{and an isomorphism}\quad
 X_0\simeq(\widetilde B_0\times F)/\pi_1(B_0)
\]
over $B_0$; the group acts by deck transformations on $\widetilde B_0$
and through $\rho_0$ on $F$.  The rational linear equivalence of
$K_{B_0}$ is the conclusion of Theorem~4.12, not merely the numerical
triviality in the introductory statement of Theorem~1.1.  The same
formulation appears in M\"uller~\cite[Theorem~2.1]{Muller2025}.
Cao--H\"oring~\cite{CaoHoring}
provides the earlier smooth locally trivial structure theorem; the
constant transition maps used here come from Matsumura--Wang.

Purity of the branch locus makes $p_0$ \'etale, since its target $X$
is smooth.  Hence $X_0$ is smooth and projective.  Local product charts
for $f_0$ imply that $B_0$ and $F$ are smooth as well: the embedding
dimension and the dimension of a product of analytic germs are sums,
so a smooth product has smooth factors.  The fiber $F$ is projective
because it is closed in $X_0$.  On the fiber over $b\in B_0$ the tangent
sequence gives
\[
 K_{X_0}^{-1}|_F\simeq K_F^{-1}\otimes
 \bigl(\det T_bB_0\bigr)\otimes\mathcal O_F.
\]
A choice of nonzero vector in the constant determinant line therefore
transfers the restricted smooth semipositive metric to $-K_F$.
Rational connectedness and Hodge symmetry give
$H^i(F,\mathcal O_F)=0$ for every $i>0$.  Thus $F$ has all the
nonequivariant hypotheses of
Theorem~\ref{V:controlled-invariant-nonvanishing}.

\subsection{Making the monodromy group connected and commutative}

The structure theorem of Demailly--Peternell--Schneider
\cite[Structure Theorem~(iii), p.~218]{DPS1996} shows that
$\pi_1(X_0)$ is virtually abelian.  Since $f_0$ is a fiber bundle with
connected fiber, its homotopy sequence makes
$\pi_1(X_0)\to\pi_1(B_0)$ surjective.  The image of an abelian
finite-index subgroup is again abelian of finite index.  Consequently
$\pi_1(B_0)$ is virtually abelian.

To place monodromy in a fixed algebraic group, choose an ample line
bundle $A$ on $X_0$ and put $A_F=A|_F$.  The restrictions of $c_1(A)$
form a flat section of $R^2(f_0)_*\mathbb Z$, so monodromy fixes the
integral class $c_1(A_F)$.  Since $H^1(F,\mathcal O_F)=0$, the exponential
sequence makes the Chern-class map
$\operatorname{Pic}(F)\to H^2(F,\mathbb Z)$ injective.  It follows that
$\rho_0(\gamma)^*A_F\simeq A_F$ for every $\gamma\in\pi_1(B_0)$.
Choose $a>0$ such that $A_F^a$ is very ample.  The projective embedding
by $A_F^a$ identifies $\operatorname{Aut}(F,A_F^a)$ with the closed stabilizer
of the embedded $F$ in $\operatorname{PGL}(H^0(F,A_F^a)^*)$.
Monodromy therefore lies in a linear algebraic group.  This argument
uses a projective action and requires no choice of linearization on
$A$ or $A_F$.

Pass to the connected finite \'etale cover of $B_0$ corresponding to
an abelian subgroup of finite index in $\pi_1(B_0)$, and pull back
$X_0$.  The resulting monodromy has commutative Zariski closure $H$:
the commutator morphism is trivial on the dense product of the
monodromy image with itself.  Pass next to the cover corresponding to
the inverse image of the identity component $H^\circ$.  The new
monodromy is dense in $H^\circ$, since each component of an algebraic
group is open and closed.  Write
\[
 f:Y\longrightarrow B,\qquad p:Y\longrightarrow X,
 \qquad G=\overline{\rho(\pi_1(B))}^{\mathrm{Zar}}=H^\circ
\]
for the resulting fibration, composite finite \'etale cover, and
monodromy closure.  These finite covers of projective varieties are
projective, and $K_B$ is still torsion.  In particular $G$ is a
connected commutative linear algebraic group acting on the same
fiber $F$.

\subsection{From a torus-fixed section to a section on \texorpdfstring{$X$}{X}}

We use the natural action on anticanonical bundles: an automorphism
acts on $K_F^{-1}=\det T_F$ by the determinant of its differential.
It induces an algebraic representation on every finite-dimensional
space $H^0(F,K_F^{-m})$.  This specified action is also the action in
the relative canonical transition functions of the locally constant
fibration.

Let $T\subset G$ be a maximal algebraic torus.  Theorem~\ref{V:controlled-invariant-nonvanishing}
gives an integer $m>0$ for which
\[
 W:=H^0(F,K_F^{-m})^T\ne0.
\]
If $F$ is a point, take $m=1$ and the constant section; the rest of the
argument is unchanged.  The trivial torus is allowed.  The inclusion
$G\subset\operatorname{Aut}(F)$ makes the action faithful, so there
is no ineffective subgroup that could introduce a scalar character
on the natural anticanonical line.

For completeness, the following elementary argument explains why
$W$ contains a vector fixed by all of $G$.  Jordan decomposition in a
connected commutative linear algebraic group over $\mathbb C$ gives
\[
 G=T\times U,
\]
where $U$ is connected and unipotent.  Commutativity makes $W$ stable
under $U$.  Its Lie algebra acts on $W$ by commuting nilpotent
endomorphisms $N_1,\ldots,N_e$, after choosing a basis of the Lie
algebra.  Their kernels have nonzero intersection: successively
restrict $N_j$ to the intersection of the preceding kernels, which
it preserves by commutativity; a nilpotent endomorphism of a nonzero
finite-dimensional space has nonzero kernel.  A nonzero vector in
this intersection is fixed by $U$, because $U$ is connected and is
generated by the exponentials of its Lie algebra.  We obtain
\[
 0\ne s\in W^U=H^0(F,K_F^{-m})^G.
\]
This is the connected-group step in
\cite[Lemma~2.3 and the proof of Theorem~2.2]{Muller2025}; we have
retained the unipotent argument because $G$ need not be a torus.
The input here is a fixed vector for the natural linearization, not
a section of an unspecified character.

Since $s$ is fixed by monodromy, its constant extension to
$\widetilde B\times F$ descends to a nonzero section
\[
 \sigma\in H^0(Y,K_{Y/B}^{-m}).
\]
The canonical bundle formula reads
$K_Y=K_{Y/B}\otimes f^*K_B$.  Choose $r>0$ with
$K_B^r\simeq\mathcal O_B$ and a nowhere-zero section
$\tau\in H^0(B,K_B^{-r})$.  Then
\[
 \sigma^r\otimes f^*\tau^m
 \in H^0(Y,K_Y^{-mr})\setminus\{0\}.
\]
All these holomorphic sections are algebraic by projectivity.  Finally,
let $d=\deg p$.  Since $p$ is finite \'etale,
$K_Y^{-mr}\simeq p^*K_X^{-mr}$, and the norm of this section is a
section of $K_X^{-mrd}$.  It is nonzero: over the generic point the
norm is a product of nonzero conjugates in a finite field extension.
Thus $X$ has a nonzero section of a positive anticanonical power.
The construction uses the locally constant product only after passing
to the universal cover of the base; no product decomposition of a
finite cover of $X$ is required.
